\documentclass[11pt,a4paper]{article}
\usepackage[utf8]{inputenc}
\usepackage[T1]{fontenc}
\usepackage[spanish,es-minimal,es-noshorthands]{babel}
\AtBeginDocument{\renewcommand{\sin}{\operatorname{sen}}\renewcommand{\arcsin}{\operatorname{arcsen}}}
\usepackage{lmodern}
\usepackage[sfdefault,lining]{FiraSans}
\usepackage[a4paper,top=2.3cm,bottom=2.2cm,left=1.8cm,right=1.8cm,headheight=1.3cm,headsep=0.4cm,footskip=1cm]{geometry}
\usepackage{amsmath,amssymb}
\usepackage{icomma}
\usepackage{graphicx}
\usepackage[table]{xcolor}
\usepackage{tikz}
\usepackage{fontawesome5}
\usepackage{ccicons}
\usepackage{fancyhdr}
\usepackage{lastpage}
\usepackage{tasks}
\usepackage{needspace}
\usepackage[most]{tcolorbox}
\usepackage{array}
\usepackage[hidelinks]{hyperref}
\AtBeginDocument{\hypersetup{pdftitle={Ficha de ejercicios de la unidad 8: Integrales. 2º Bach.},pdfsubject={Matemáticas CCSS 2º Bach.},pdfkeywords={integrales, primitiva, integral indefinida, integrales inmediatas, integral definida, regla de Barrow, áreas, área entre dos curvas}}}

\definecolor{navy}{HTML}{1B2748}
\definecolor{accent}{HTML}{E8603A}
\definecolor{teal}{HTML}{1F9E90}
\definecolor{soft}{HTML}{EEF1F8}
\definecolor{muted}{HTML}{5D6685}

\setlength{\parindent}{0pt}
\setlength{\parskip}{3pt}
\settasks{label=\alph*), label-format=\bfseries\color{navy}, label-width=1.3em, item-indent=1.8em, column-sep=0.6em, after-item-skip=5pt, before-skip=3pt, after-skip=4pt}

% ---------- capçalera i peu ----------
\newcommand{\llicencia}{\href{https://creativecommons.org/licenses/by-nc-sa/4.0/deed.es}{\ccbyncsa}\ \ Licencia Creative Commons Atribución-NoComercial-CompartirIgual 4.0 Internacional (CC BY-NC-SA 4.0)}
\pagestyle{fancy}
\fancyhf{}
\fancyhead[L]{\raisebox{-0.35\height}{\includegraphics[height=0.95cm]{logo-gva.png}}}
\fancyhead[C]{\small\color{muted}Ficha de ejercicios · Unidad 8: Integrales · 2\textsuperscript{o} Bach.}
\fancyhead[R]{\raisebox{-0.35\height}{\includegraphics[height=0.95cm]{logo-vermellar.png}}}
\fancyfoot[L]{\scriptsize\color{muted}\llicencia}
\fancyfoot[R]{\small\color{muted}\thepage}
\renewcommand{\headrulewidth}{0.4pt}
\renewcommand{\headrule}{\color{muted!40}\hrule height \headrulewidth}
\fancypagestyle{primera}{\fancyhf{}\renewcommand{\headrulewidth}{0pt}
  \fancyfoot[L]{\scriptsize\color{muted}\llicencia}
  \fancyfoot[R]{\small\color{muted}\thepage}}

% ---------- nivells i exercicis ----------
\newcommand{\estrelles}[1]{{\color{accent}\ifcase#1\or\faStar\faStar[regular]\faStar[regular]\or\faStar\faStar\faStar[regular]\or\faStar\faStar\faStar\fi}}
\newcounter{exer}
\newcommand{\ex}[1]{\par\needspace{4\baselineskip}\medskip\refstepcounter{exer}\noindent{\color{navy}\bfseries\theexer.}\ \ }

\newcommand{\seccio}[2]{\par\needspace{10\baselineskip}\bigskip
  \begin{tcolorbox}[colback=navy,colframe=navy,coltext=white,arc=2mm,boxrule=0pt,left=3mm,right=3mm,top=1.2mm,bottom=1.2mm]
  {\large\bfseries #1\quad #2}
  \end{tcolorbox}\vspace{-2pt}}


% quadre de recordatori (només a la fitxa d'exercicis)
\newcommand{\recorda}[1]{\par\smallskip\begin{tcolorbox}[colback=accent!7,colframe=accent!55,arc=2mm,boxrule=0.5pt,left=2.5mm,right=2.5mm,top=1mm,bottom=1mm]\small{\color{accent}\faLightbulb[regular]}\ \textbf{Recuerda:}\ #1\end{tcolorbox}}
\newcommand{\res}[1]{#1}\newcommand{\rest}[1]{#1}

% =====================================================================
% Capçalera FentMates de les fitxes: un color per a cada curs.
% Els colors són els de la pàgina inicial (LEVELS a assets/js/materials.js);
% el logo ve de _eines/logo/pdf/ (python3 _eines/logo/genera.py).
%
% Ús, al final del preàmbul de la unitat:
%   % =====================================================================
% Capçalera FentMates de les fitxes: un color per a cada curs.
% Els colors són els de la pàgina inicial (LEVELS a assets/js/materials.js);
% el logo ve de _eines/logo/pdf/ (python3 _eines/logo/genera.py).
%
% Ús, al final del preàmbul de la unitat:
%   % =====================================================================
% Capçalera FentMates de les fitxes: un color per a cada curs.
% Els colors són els de la pàgina inicial (LEVELS a assets/js/materials.js);
% el logo ve de _eines/logo/pdf/ (python3 _eines/logo/genera.py).
%
% Ús, al final del preàmbul de la unitat:
%   \input{../comu/fentmates.tex}
%   \fmcurs{1eso}
% i, al començament del document:
%   \capcalera{Matemàtiques}{1r ESO}{Fitxa d'exercicis · Unitat 2}{Divisibilitat}
% Opcional abans de \input: \def\fmneed{14} (espai mínim abans de cada secció).
% =====================================================================
\makeatletter
\def\fm@logodir{../../logo/pdf}
\providecommand{\fmneed}{10}
\def\fm@cursnom{}

% clau del curs, color i color del text damunt del color
\newcommand{\fm@curs}[3]{\@namedef{fm@color@#1}{#2}\@namedef{fm@text@#1}{#3}}
\fm@curs{1eso}{E6B84A}{navy}
\fm@curs{2eso}{3B82F6}{white}
\fm@curs{3eso}{FF7A50}{white}
\fm@curs{4eso-a}{A78BFA}{white}
\fm@curs{4eso-b}{6D28D9}{white}
\fm@curs{1batx-mi}{9B1C1C}{white}
\fm@curs{1batx-ccss}{EF6B6B}{white}
\fm@curs{2batx-mii}{1E6B47}{white}
\fm@curs{2batx-ccss}{2BB3A3}{white}

\newcommand{\fmcurs}[1]{%
  \@ifundefined{fm@color@#1}{\PackageError{fentmates}{Curs desconegut: #1}{}}{}%
  \edef\fm@tmp{\noexpand\definecolor{curs}{HTML}{\@nameuse{fm@color@#1}}%
               \noexpand\colorlet{curstext}{\@nameuse{fm@text@#1}}}\fm@tmp
  \colorlet{cursfosc}{curs!72!black}%
  \def\fm@logo{\fm@logodir/vertical-#1.pdf}%
  \def\fm@marca{\fm@logodir/marca-#1.pdf}}

% ---------- capçalera de la primera pàgina ----------
% #1 matèria (petit), #2 curs, #3 tipus de document i unitat, #4 títol
\newcommand{\capcalera}[4]{%
  \gdef\fm@cursnom{#2}%
  \par\noindent
  \begin{tcolorbox}[colback=curs,colframe=curs,boxrule=0pt,arc=2.2mm,width=3.5cm,height=2.15cm,
      valign=center,left=3mm,right=2mm,top=1mm,bottom=1mm,nobeforeafter,box align=center]
    \color{curstext}{\footnotesize #1}\par\vspace{1pt}{\fontsize{19}{21}\selectfont\bfseries #2\par}%
  \end{tcolorbox}\hspace{5mm}%
  \begin{minipage}[c]{\dimexpr\linewidth-3.5cm-5mm-2.9cm\relax}
    {\color{cursfosc}\small\bfseries\MakeUppercase{#3}}\par\vspace{2pt}
    {\color{navy}\fontsize{20}{23}\selectfont\bfseries #4\par}
  \end{minipage}\hfill
  \begin{minipage}[c]{2.6cm}\raggedleft\includegraphics[width=2.6cm]{\fm@logo}\end{minipage}\par
  \vspace{2pt}}

% ---------- barres de secció del color del curs ----------
\renewcommand{\seccio}[2]{\par\needspace{\fmneed\baselineskip}\bigskip
  \begin{tcolorbox}[colback=curs,colframe=curs,coltext=curstext,arc=2mm,boxrule=0pt,left=3mm,right=3mm,top=1.2mm,bottom=1.2mm]
  {\large\bfseries #1\quad #2}
  \end{tcolorbox}\vspace{-2pt}}

% ---------- estrelles de dificultat del color del curs (a Batxillerat no se'n posen) ----------
\@ifundefined{estrelles}{}{%
  \renewcommand{\estrelles}[1]{{\color{curs}\ifcase#1\or\faStar\faStar[regular]\faStar[regular]\or\faStar\faStar\faStar[regular]\or\faStar\faStar\faStar\fi}}}

% ---------- capçalera i peu de les pàgines ----------
% primera pàgina: logos institucionals al peu
\fancypagestyle{primera}{\fancyhf{}\renewcommand{\headrulewidth}{0pt}
  \fancyfoot[L]{\raisebox{-0.32\height}{\includegraphics[height=0.78cm]{logo-gva.png}}\hspace{5mm}%
    \raisebox{-0.32\height}{\includegraphics[height=0.78cm]{logo-vermellar.png}}}
  \fancyfoot[R]{\scriptsize\color{muted}\href{https://creativecommons.org/licenses/by-nc-sa/4.0/deed.ca}{\ccbyncsa}\ \ CC BY-NC-SA 4.0 · fentmates.org\qquad\small\thepage}}
% la resta: el símbol a l'esquerra i el curs a la dreta, amb el títol al centre com abans
\fancyhead[L]{\raisebox{-0.3\height}{\includegraphics[height=0.62cm]{\fm@marca}}}
\fancyhead[R]{\small\bfseries\color{cursfosc}\fm@cursnom}
\renewcommand{\headrule}{\color{curs!45}\hrule height \headrulewidth}
\makeatother

%   \fmcurs{1eso}
% i, al començament del document:
%   \capcalera{Matemàtiques}{1r ESO}{Fitxa d'exercicis · Unitat 2}{Divisibilitat}
% Opcional abans de \input: \def\fmneed{14} (espai mínim abans de cada secció).
% =====================================================================
\makeatletter
\def\fm@logodir{../../logo/pdf}
\providecommand{\fmneed}{10}
\def\fm@cursnom{}

% clau del curs, color i color del text damunt del color
\newcommand{\fm@curs}[3]{\@namedef{fm@color@#1}{#2}\@namedef{fm@text@#1}{#3}}
\fm@curs{1eso}{E6B84A}{navy}
\fm@curs{2eso}{3B82F6}{white}
\fm@curs{3eso}{FF7A50}{white}
\fm@curs{4eso-a}{A78BFA}{white}
\fm@curs{4eso-b}{6D28D9}{white}
\fm@curs{1batx-mi}{9B1C1C}{white}
\fm@curs{1batx-ccss}{EF6B6B}{white}
\fm@curs{2batx-mii}{1E6B47}{white}
\fm@curs{2batx-ccss}{2BB3A3}{white}

\newcommand{\fmcurs}[1]{%
  \@ifundefined{fm@color@#1}{\PackageError{fentmates}{Curs desconegut: #1}{}}{}%
  \edef\fm@tmp{\noexpand\definecolor{curs}{HTML}{\@nameuse{fm@color@#1}}%
               \noexpand\colorlet{curstext}{\@nameuse{fm@text@#1}}}\fm@tmp
  \colorlet{cursfosc}{curs!72!black}%
  \def\fm@logo{\fm@logodir/vertical-#1.pdf}%
  \def\fm@marca{\fm@logodir/marca-#1.pdf}}

% ---------- capçalera de la primera pàgina ----------
% #1 matèria (petit), #2 curs, #3 tipus de document i unitat, #4 títol
\newcommand{\capcalera}[4]{%
  \gdef\fm@cursnom{#2}%
  \par\noindent
  \begin{tcolorbox}[colback=curs,colframe=curs,boxrule=0pt,arc=2.2mm,width=3.5cm,height=2.15cm,
      valign=center,left=3mm,right=2mm,top=1mm,bottom=1mm,nobeforeafter,box align=center]
    \color{curstext}{\footnotesize #1}\par\vspace{1pt}{\fontsize{19}{21}\selectfont\bfseries #2\par}%
  \end{tcolorbox}\hspace{5mm}%
  \begin{minipage}[c]{\dimexpr\linewidth-3.5cm-5mm-2.9cm\relax}
    {\color{cursfosc}\small\bfseries\MakeUppercase{#3}}\par\vspace{2pt}
    {\color{navy}\fontsize{20}{23}\selectfont\bfseries #4\par}
  \end{minipage}\hfill
  \begin{minipage}[c]{2.6cm}\raggedleft\includegraphics[width=2.6cm]{\fm@logo}\end{minipage}\par
  \vspace{2pt}}

% ---------- barres de secció del color del curs ----------
\renewcommand{\seccio}[2]{\par\needspace{\fmneed\baselineskip}\bigskip
  \begin{tcolorbox}[colback=curs,colframe=curs,coltext=curstext,arc=2mm,boxrule=0pt,left=3mm,right=3mm,top=1.2mm,bottom=1.2mm]
  {\large\bfseries #1\quad #2}
  \end{tcolorbox}\vspace{-2pt}}

% ---------- estrelles de dificultat del color del curs (a Batxillerat no se'n posen) ----------
\@ifundefined{estrelles}{}{%
  \renewcommand{\estrelles}[1]{{\color{curs}\ifcase#1\or\faStar\faStar[regular]\faStar[regular]\or\faStar\faStar\faStar[regular]\or\faStar\faStar\faStar\fi}}}

% ---------- capçalera i peu de les pàgines ----------
% primera pàgina: logos institucionals al peu
\fancypagestyle{primera}{\fancyhf{}\renewcommand{\headrulewidth}{0pt}
  \fancyfoot[L]{\raisebox{-0.32\height}{\includegraphics[height=0.78cm]{logo-gva.png}}\hspace{5mm}%
    \raisebox{-0.32\height}{\includegraphics[height=0.78cm]{logo-vermellar.png}}}
  \fancyfoot[R]{\scriptsize\color{muted}\href{https://creativecommons.org/licenses/by-nc-sa/4.0/deed.ca}{\ccbyncsa}\ \ CC BY-NC-SA 4.0 · fentmates.org\qquad\small\thepage}}
% la resta: el símbol a l'esquerra i el curs a la dreta, amb el títol al centre com abans
\fancyhead[L]{\raisebox{-0.3\height}{\includegraphics[height=0.62cm]{\fm@marca}}}
\fancyhead[R]{\small\bfseries\color{cursfosc}\fm@cursnom}
\renewcommand{\headrule}{\color{curs!45}\hrule height \headrulewidth}
\makeatother

%   \fmcurs{1eso}
% i, al començament del document:
%   \capcalera{Matemàtiques}{1r ESO}{Fitxa d'exercicis · Unitat 2}{Divisibilitat}
% Opcional abans de \input: \def\fmneed{14} (espai mínim abans de cada secció).
% =====================================================================
\makeatletter
\def\fm@logodir{../../logo/pdf}
\providecommand{\fmneed}{10}
\def\fm@cursnom{}

% clau del curs, color i color del text damunt del color
\newcommand{\fm@curs}[3]{\@namedef{fm@color@#1}{#2}\@namedef{fm@text@#1}{#3}}
\fm@curs{1eso}{E6B84A}{navy}
\fm@curs{2eso}{3B82F6}{white}
\fm@curs{3eso}{FF7A50}{white}
\fm@curs{4eso-a}{A78BFA}{white}
\fm@curs{4eso-b}{6D28D9}{white}
\fm@curs{1batx-mi}{9B1C1C}{white}
\fm@curs{1batx-ccss}{EF6B6B}{white}
\fm@curs{2batx-mii}{1E6B47}{white}
\fm@curs{2batx-ccss}{2BB3A3}{white}

\newcommand{\fmcurs}[1]{%
  \@ifundefined{fm@color@#1}{\PackageError{fentmates}{Curs desconegut: #1}{}}{}%
  \edef\fm@tmp{\noexpand\definecolor{curs}{HTML}{\@nameuse{fm@color@#1}}%
               \noexpand\colorlet{curstext}{\@nameuse{fm@text@#1}}}\fm@tmp
  \colorlet{cursfosc}{curs!72!black}%
  \def\fm@logo{\fm@logodir/vertical-#1.pdf}%
  \def\fm@marca{\fm@logodir/marca-#1.pdf}}

% ---------- capçalera de la primera pàgina ----------
% #1 matèria (petit), #2 curs, #3 tipus de document i unitat, #4 títol
\newcommand{\capcalera}[4]{%
  \gdef\fm@cursnom{#2}%
  \par\noindent
  \begin{tcolorbox}[colback=curs,colframe=curs,boxrule=0pt,arc=2.2mm,width=3.5cm,height=2.15cm,
      valign=center,left=3mm,right=2mm,top=1mm,bottom=1mm,nobeforeafter,box align=center]
    \color{curstext}{\footnotesize #1}\par\vspace{1pt}{\fontsize{19}{21}\selectfont\bfseries #2\par}%
  \end{tcolorbox}\hspace{5mm}%
  \begin{minipage}[c]{\dimexpr\linewidth-3.5cm-5mm-2.9cm\relax}
    {\color{cursfosc}\small\bfseries\MakeUppercase{#3}}\par\vspace{2pt}
    {\color{navy}\fontsize{20}{23}\selectfont\bfseries #4\par}
  \end{minipage}\hfill
  \begin{minipage}[c]{2.6cm}\raggedleft\includegraphics[width=2.6cm]{\fm@logo}\end{minipage}\par
  \vspace{2pt}}

% ---------- barres de secció del color del curs ----------
\renewcommand{\seccio}[2]{\par\needspace{\fmneed\baselineskip}\bigskip
  \begin{tcolorbox}[colback=curs,colframe=curs,coltext=curstext,arc=2mm,boxrule=0pt,left=3mm,right=3mm,top=1.2mm,bottom=1.2mm]
  {\large\bfseries #1\quad #2}
  \end{tcolorbox}\vspace{-2pt}}

% ---------- estrelles de dificultat del color del curs (a Batxillerat no se'n posen) ----------
\@ifundefined{estrelles}{}{%
  \renewcommand{\estrelles}[1]{{\color{curs}\ifcase#1\or\faStar\faStar[regular]\faStar[regular]\or\faStar\faStar\faStar[regular]\or\faStar\faStar\faStar\fi}}}

% ---------- capçalera i peu de les pàgines ----------
% primera pàgina: logos institucionals al peu
\fancypagestyle{primera}{\fancyhf{}\renewcommand{\headrulewidth}{0pt}
  \fancyfoot[L]{\raisebox{-0.32\height}{\includegraphics[height=0.78cm]{logo-gva.png}}\hspace{5mm}%
    \raisebox{-0.32\height}{\includegraphics[height=0.78cm]{logo-vermellar.png}}}
  \fancyfoot[R]{\scriptsize\color{muted}\href{https://creativecommons.org/licenses/by-nc-sa/4.0/deed.ca}{\ccbyncsa}\ \ CC BY-NC-SA 4.0 · fentmates.org\qquad\small\thepage}}
% la resta: el símbol a l'esquerra i el curs a la dreta, amb el títol al centre com abans
\fancyhead[L]{\raisebox{-0.3\height}{\includegraphics[height=0.62cm]{\fm@marca}}}
\fancyhead[R]{\small\bfseries\color{cursfosc}\fm@cursnom}
\renewcommand{\headrule}{\color{curs!45}\hrule height \headrulewidth}
\makeatother

\fmcurs{2batx-ccss}

\begin{document}
\thispagestyle{primera}

\capcalera{Matemáticas CCSS}{2\textsuperscript{o} Bach.}{Ficha de ejercicios · Unidad 8}{\raggedright Integrales}
\vspace{10pt}
Nombre y apellidos: \rule{6.6cm}{0.4pt}\hfill Grupo: \rule{1.6cm}{0.4pt}\hfill Fecha: \rule{2.2cm}{0.4pt}
\vspace{8pt}

\begin{tcolorbox}[colback=soft,colframe=soft,arc=2mm,left=3mm,right=3mm,top=1.5mm,bottom=1.5mm]
\small Resuelve los ejercicios en el cuaderno mostrando todos los pasos. En la PAU de Ciencias Sociales las integrales son de polinomios: integra término a término, aplica la regla de Barrow con cuidado con los signos y, en los problemas de áreas, busca siempre primero los puntos de corte y haz un esbozo del recinto. No olvides la constante $C$ en las integrales indefinidas ni las unidades ($\text{u}^2$) en las áreas.
\end{tcolorbox}

% =====================================================
\seccio{1}{Primitivas e integral indefinida}

\ex{1} Comprueba, derivando, si $F$ es una primitiva de $f$ en cada caso:\begin{tasks}(2)
  \task $F(x)=x^3-5x+2$, \ $f(x)=3x^2-5$
  \task $F(x)=(x^2-1)^2$, \ $f(x)=4x^3-4x$
  \task $F(x)=\dfrac{x^2}{2}+\dfrac1x$, \ $f(x)=x+\dfrac{1}{x^2}$
  \task $F(x)=\dfrac{(2x+3)^3}{6}$, \ $f(x)=(2x+3)^2$
\end{tasks}\recorda{$F$ es una primitiva de $f$ si $F'(x)=f(x)$.}

\ex{1} Escribe tres primitivas diferentes de $f(x)=3x^2-4$. ¿Qué relación hay entre sus gráficas?

\ex{2} Di si las afirmaciones siguientes son verdaderas o falsas. Justifícalo.\begin{tasks}(1)
  \task Si $F$ y $G$ son dos primitivas de la misma función $f$, entonces $F-G$ es constante.
  \task $\displaystyle\int f(x)\cdot g(x)\,dx=\int f(x)\,dx\cdot\int g(x)\,dx$.
  \task $\displaystyle\int 5f(x)\,dx=5\int f(x)\,dx$.
\end{tasks}


% =====================================================
\seccio{2}{Integrales inmediatas}

\ex{1} Calcula las integrales de polinomios siguientes:\begin{tasks}(3)
  \task $\displaystyle\int (4x^3-6x^2+2x-5)\,dx$
  \task $\displaystyle\int (x^2+3x-1)\,dx$
  \task $\displaystyle\int (5x^4-x)\,dx$
  \task $\displaystyle\int \left(2-\frac x3\right)\,dx$
  \task $\displaystyle\int \left(\frac{x^3}{2}-\frac{2x^2}{3}+1\right)\,dx$
  \task $\displaystyle\int (6x^5+4x^3)\,dx$
  \task $\displaystyle\int 7\,dx$
  \task $\displaystyle\int (0{,}5x-1{,}2)\,dx$
  \task $\displaystyle\int (-3x^2+8x)\,dx$
\end{tasks}\recorda{Integramos término a término: $\displaystyle\int a\,x^n\,dx=a\,\frac{x^{n+1}}{n+1}+C$ y $\displaystyle\int k\,dx=kx+C$. No olvides la constante $C$.}

\ex{1} Desarrolla primero y después integra:\begin{tasks}(3)
  \task $\displaystyle\int x(x+4)\,dx$
  \task $\displaystyle\int (x-3)^2\,dx$
  \task $\displaystyle\int (x+1)(x-1)\,dx$
  \task $\displaystyle\int (2x-1)(x+2)\,dx$
  \task $\displaystyle\int x^2(3x-2)\,dx$
  \task $\displaystyle\int \frac{x^3-2x^2}{x}\,dx$
\end{tasks}\recorda{No hay ninguna regla para integrar un producto o un cociente: si son polinomios, primero los desarrollamos (o simplificamos).}

\ex{2} Calcula estas integrales inmediatas (escribe las raíces y las fracciones como potencias):\begin{tasks}(2)
  \task $\displaystyle\int \frac{1}{x^3}\,dx$
  \task $\displaystyle\int \sqrt[3]{x}\,dx$
  \task $\displaystyle\int \left(\frac2x+3e^x\right)\,dx$
  \task $\displaystyle\int \left(x+\frac{1}{\sqrt x}\right)\,dx$
\end{tasks}\recorda{$\displaystyle\int x^n\,dx=\frac{x^{n+1}}{n+1}+C$ también vale para $n$ negativo o fraccionario (si $n\neq-1$). Además, $\displaystyle\int\frac1x\,dx=\ln|x|+C$ y $\displaystyle\int e^x\,dx=e^x+C$.}


% =====================================================
\seccio{3}{Integrales casi inmediatas}

\ex{2} Calcula las integrales siguientes ajustando constantes:\begin{tasks}(2)
  \task $\displaystyle\int (5x+2)^3\,dx$
  \task $\displaystyle\int 3x^2(x^3-1)^4\,dx$
  \task $\displaystyle\int x(x^2+5)^2\,dx$
  \task $\displaystyle\int \frac{4x+2}{x^2+x+1}\,dx$
  \task $\displaystyle\int \frac{1}{2x-5}\,dx$
  \task $\displaystyle\int e^{-2x}\,dx$
  \task $\displaystyle\int (2x+1)\,e^{x^2+x}\,dx$
  \task $\displaystyle\int x^3e^{x^4}\,dx$
\end{tasks}\recorda{$\displaystyle\int u'\,u^n\,dx=\frac{u^{n+1}}{n+1}+C$, \ $\displaystyle\int\frac{u'}{u}\,dx=\ln|u|+C$, \ $\displaystyle\int u'\,e^{u}\,dx=e^{u}+C$. Si falta un factor constante, multiplicamos y dividimos por él.}

\ex{2} Calcula $\displaystyle\int(x+2)^2\,dx$ de dos maneras: como integral casi inmediata y desarrollando el cuadrado. ¿Sale el mismo resultado? Explícalo.


% =====================================================
\seccio{4}{Primitivas que cumplen una condición}

\ex{1} Halla la primitiva de $f(x)=6x^2-4x+1$ que cumple $F(1)=3$.

\ex{1} Halla la primitiva de $f(x)=x^2-1$ cuya gráfica pasa por el punto $(3,\,4)$.

\ex{2} La pendiente de la recta tangente a la gráfica de una función $f$ en cada punto $x$ es $3x^2-2x$, y la gráfica pasa por el punto $(1,\,2)$. Halla $f$.

\ex{2} De una función $f$ sabemos que $f''(x)=6x-2$, que $f'(1)=3$ y que $f(0)=1$. Halla $f$.

\ex{2} El coste marginal de fabricar $x$ unidades de un producto es $C'(x)=0{,}6x+20$ euros por unidad, y los costes fijos son de 500 €.\begin{tasks}(1)
  \task Halla la función de coste total $C(x)$.
  \task ¿Cuánto cuesta fabricar 50 unidades?
\end{tasks}\recorda{El coste total es una primitiva del coste marginal. Los costes fijos son el coste cuando no se fabrica nada: $C(0)$.}

\ex{2} La velocidad de crecimiento de una población, en miles de habitantes por año, es $P'(t)=0{,}3t^2+2t$, donde $t$ son los años transcurridos desde el año 2020. En 2020 la población era de 50 mil habitantes.\begin{tasks}(1)
  \task Halla la función $P(t)$ que da la población.
  \task ¿Qué población habrá en 2025?
\end{tasks}


% =====================================================
\seccio{5}{La integral definida. Regla de Barrow}

\ex{1} Calcula estas integrales definidas haciendo un dibujo y usando áreas de figuras conocidas:\begin{tasks}(4)
  \task $\displaystyle\int_{0}^{5} 4\,dx$
  \task $\displaystyle\int_{-2}^{2} (x+2)\,dx$
  \task $\displaystyle\int_{0}^{3} (-2x)\,dx$
  \task $\displaystyle\int_{-1}^{3} x\,dx$
\end{tasks}\recorda{Si $f\geq0$, $\int_a^b f$ es el área; las partes por debajo del eje cuentan con signo negativo.}

\ex{1} Calcula aplicando la regla de Barrow:\begin{tasks}(3)
  \task $\displaystyle\int_{0}^{3} (x^2+1)\,dx$
  \task $\displaystyle\int_{1}^{2} (6x^2-2x)\,dx$
  \task $\displaystyle\int_{-1}^{2} (3x^2-2x+1)\,dx$
  \task $\displaystyle\int_{-2}^{0} (x^3+2x)\,dx$
  \task $\displaystyle\int_{0}^{1} (x^4-x)\,dx$
  \task $\displaystyle\int_{1}^{4} \left(2-\frac x2\right)\,dx$
  \task $\displaystyle\int_{-1}^{1} (x^5-x^3+4)\,dx$
  \task $\displaystyle\int_{0}^{2} (x-1)(x+3)\,dx$
  \task $\displaystyle\int_{1}^{2} (x+1)^2\,dx$
\end{tasks}\recorda{$\displaystyle\int_a^b f(x)\,dx=\Big[F(x)\Big]_a^b=F(b)-F(a)$. Cuidado con los paréntesis cuando $F(a)$ es negativo.}

\ex{2} Calcula:\begin{tasks}(4)
  \task $\displaystyle\int_{1}^{e} \frac2x\,dx$
  \task $\displaystyle\int_{0}^{\ln2} e^x\,dx$
  \task $\displaystyle\int_{1}^{4} \frac{1}{\sqrt x}\,dx$
  \task $\displaystyle\int_{1}^{2} \frac{1}{x^2}\,dx$
\end{tasks}

\ex{1} Sabemos que $\displaystyle\int_0^3 f(x)\,dx=4$, $\displaystyle\int_3^5 f(x)\,dx=-1$ y $\displaystyle\int_0^5 g(x)\,dx=2$. Calcula:\begin{tasks}(3)
  \task $\displaystyle\int_{0}^{5} f(x)\,dx$
  \task $\displaystyle\int_{5}^{0} \big(2f(x)-g(x)\big)\,dx$
  \task $\displaystyle\int_{3}^{3} f(x)\,dx$
\end{tasks}

\ex{2} Calcula el valor de $a>0$ para que $\displaystyle\int_0^a(2x+1)\,dx=6$.

\ex{2} Calcula $k$ para que $\displaystyle\int_1^3(kx-1)\,dx=10$.

\ex{3} Halla $a$ y $b$ para que la función $f(x)=ax^2+b$ cumpla $f(1)=2$ y $\displaystyle\int_0^1 f(x)\,dx=1$.

\ex{2} Aproxima $\displaystyle\int_0^2x^2\,dx$ con 4 rectángulos de la misma anchura, tomando como altura el valor de la función en el extremo derecho de cada trozo. Compáralo con el valor exacto.


% =====================================================
\seccio{6}{Área entre una curva y el eje OX}

\ex{2} Calcula el área del recinto limitado por la gráfica de cada función y el eje OX. Haz un esbozo.\begin{tasks}(3)
  \task $f(x)=x^2-9$
  \task $f(x)=-x^2+2x+3$
  \task $f(x)=x^2-5x+4$
  \task $f(x)=x^3-9x$
  \task $f(x)=x^3-3x^2+2x$
  \task $f(x)=x^4-x^2$
\end{tasks}\recorda{Primero buscamos los cortes con el eje OX ($f(x)=0$); son los límites de integración. Si hay más de dos, integramos cada trozo por separado y sumamos los valores absolutos.}

\ex{2} Calcula el área de la región limitada por la gráfica de $f$, el eje OX y las rectas indicadas:\begin{tasks}(2)
  \task $f(x)=x^2+1$, \ $x=0$, \ $x=3$
  \task $f(x)=4-x^2$, \ $x=0$, \ $x=3$
  \task $f(x)=x^2-4x$, \ $x=-1$, \ $x=1$
  \task $f(x)=x^3-x$, \ $x=0$, \ $x=2$
  \task $f(x)=x^2-2x-3$, \ $x=0$, \ $x=4$
  \task $f(x)=2x-4$, \ $x=0$, \ $x=5$
\end{tasks}\recorda{Ahora los límites son $a$ y $b$, pero debemos mirar si $f$ corta el eje OX \textbf{dentro} del intervalo: esos cortes dividen la integral en trozos. Los cortes de fuera no cuentan.}

\ex{1} Un alumno dice que el área limitada por $f(x)=x^3$, el eje OX y las rectas $x=-2$ y $x=2$ es 0, porque $\displaystyle\int_{-2}^{2}x^3\,dx=0$. ¿Tiene razón? Calcula el área correctamente.

\ex{2} Considera la función $f(x)=x^2-4x+3$.\begin{tasks}(1)
  \task Represéntala gráficamente.
  \task Calcula el área del recinto limitado por la gráfica de $f$ y el eje OX.
  \task Calcula el área del recinto limitado por la gráfica de $f$ y los dos ejes de coordenadas.
\end{tasks}

\ex{3} Calcula $a>0$ para que el área del recinto limitado por la parábola $f(x)=a^2-x^2$ y el eje OX sea 36.

\ex{2} Calcula $k>0$ para que el área de la región limitada por $f(x)=x^2$, el eje OX y la recta $x=k$ sea 9.


% =====================================================
\seccio{7}{Área entre dos curvas}

\ex{2} Calcula el área del recinto cerrado por las gráficas de las funciones:\begin{tasks}(2)
  \task $f(x)=x^2$ \ y \ $g(x)=2x$
  \task $f(x)=x^2-2$ \ y \ $g(x)=x$
  \task $f(x)=-x^2+6x$ \ y \ $g(x)=x^2-2x$
  \task $f(x)=x^3$ \ y \ $g(x)=x^2$
  \task $f(x)=x^3-2x$ \ y \ $g(x)=2x$
  \task $f(x)=4-x^2$ \ y \ $g(x)=x^2-4$
\end{tasks}\recorda{Resolvemos $f(x)=g(x)$ para hallar los puntos de corte, que son los límites. Integramos $f-g$ en cada trozo y sumamos los valores absolutos (o integramos «la de arriba menos la de abajo»).}

\ex{2} Calcula el área de la región limitada por las gráficas de $f$ y $g$ y las rectas indicadas:\begin{tasks}(2)
  \task $f(x)=x^2+1$, \ $g(x)=-x^2-1$, \ $x=-1$, \ $x=2$
  \task $f(x)=x^2$, \ $g(x)=x$, \ $x=-1$, \ $x=1$
  \task $f(x)=x^2-4$, \ $g(x)=-x+2$, \ $x=0$, \ $x=3$
  \task $f(x)=x^3$, \ $g(x)=x$, \ $x=0$, \ $x=2$
\end{tasks}

\ex{3} Calcula el área de la región limitada por la parábola $y=x^2$, la recta tangente a la parábola en $x=1$ y el eje OY.

\ex{2} Considera las funciones $f(x)=-x^2+5x$ y $g(x)=x$.\begin{tasks}(1)
  \task Represéntalas en los mismos ejes.
  \task Calcula el área del recinto cerrado por las dos gráficas.
\end{tasks}

\ex{3} Calcula $m>0$ para que el área del recinto limitado por la parábola $y=x^2$ y la recta $y=mx$ sea 36.


% =====================================================
\seccio{8}{Problemas de aplicación}

\ex{2} La función de demanda de un producto es $p(x)=100-x^2$, donde $p$ es el precio (en euros) y $x$ la cantidad demandada (en miles de unidades). El precio de mercado es de 64 €.\begin{tasks}(1)
  \task ¿Qué cantidad se demanda a este precio?
  \task El \textbf{excedente del consumidor} es el área comprendida entre la curva de demanda y la recta $p=64$, desde $x=0$ hasta la cantidad demandada. Calcúlalo.
\end{tasks}

\ex{2} El beneficio marginal de una empresa (en miles de euros) cuando vende $x$ toneladas es $B'(x)=40-2x$. Si no vende nada, tiene unas pérdidas de 100 mil euros.\begin{tasks}(1)
  \task Halla la función de beneficio $B(x)$.
  \task ¿Cuántas toneladas debe vender para obtener el beneficio máximo? ¿Cuál es ese beneficio?
  \task ¿Cuánto aumenta el beneficio si pasa de vender 5 toneladas a vender 10? (Calcúlalo con una integral definida.)
\end{tasks}

\ex{3} Un pozo de petróleo produce, $t$ años después de abrirlo, a un ritmo de $r(t)=12-\frac{3}{4}t^2$ miles de barriles por año.\begin{tasks}(1)
  \task ¿Cuándo dejará de producir?
  \task ¿Cuántos barriles producirá en total, desde que se abre hasta que se cierra?
\end{tasks}

\vfill
\begin{center}\color{muted}\small\faCheckCircle[regular]\ \ ¡Ánimo! Comprueba tus resultados con la calculadora de integrales y áreas de los apuntes interactivos.\end{center}

\end{document}
